\chapter*{ABSTRACT}
\addcontentsline{toc}{chapter}{ABSTRACT}

With the development of artificial intelligence, human-computer interaction, and mobile Internet technologies, assistive communication and training systems for hearing-impaired users have become an important application direction of accessible information technology. This thesis designs and implements HearBridge, a multi-end collaborative system for hearing-impaired communication and training. The system consists of a HarmonyOS mobile application, a Spring Boot backend, a Python-based gesture recognition service, and a Web administration platform. It supports sign language resource presentation, digital avatar playback, real-time gesture recognition, training sample collection, model training management, and model version release.

This thesis first analyzes the research background, user requirements, and functional objectives of the system. Then, it presents the overall architecture that coordinates the mobile client, backend service, recognition service, and administration platform. In the implementation, the mobile application is responsible for user interaction, camera capture, and recognition result display; the backend manages business data and API services; the Python service performs gesture recognition based on MediaPipe and deep learning models; and the Web administration platform manages sign resources, samples, training tasks, and model versions. Finally, functional tests and key process tests are conducted to verify the feasibility of the system in gesture recognition, sample collection, model release, and multi-end collaboration.

\vspace{1em}

\noindent\textbf{Keywords:} HearBridge; hearing-impaired assistance; sign language recognition; HarmonyOS; Spring Boot; deep learning
